% !TEX root = ../matan.tex

\section{Теорема о неявной функции}

\subsection*{Формулировка}

Рассмотрим отображение $F\colon W\to\RR^m$, заданное на окрестности
$W\subset\RR^d\times\RR^m$ точки $(x_0,y_0)$. Через $F'_y(x,y)$ обозначаем
производную (матрицу Якоби размера $m\times m$) по переменным $y$.

\begin{Theorem}{О неявной функции}{}
	Пусть отображение $F\colon W\to\RR^m$ удовлетворяет условиям:
	\begin{enumerate}
		\item $F(x_0,y_0)=0$;
		\item $F$ дифференцируемо в $W$, причём $F'_y$ непрерывна в точке $(x_0,y_0)$;
		\item оператор $F'_y(x_0,y_0)$ обратим (в конечномерном случае ---
		$\det F'_y(x_0,y_0)\ne 0$).
	\end{enumerate}
	Тогда существуют окрестности $U_{x_0}\subset\RR^d$ и $V_{y_0}\subset\RR^m$ и
	отображение $\vphi\colon U_{x_0}\to V_{y_0}$ такие, что в $U_{x_0}\times V_{y_0}$
	\[
	F(x,y)=0\iff y=\vphi(x),
	\]
	при этом $y_0=\vphi(x_0)$ и $\vphi$ непрерывно в точке $x_0$.
\end{Theorem}

\begin{Remark}{Нормировка}{}
	Без ограничения общности $x_0=0$, $y_0=0$.
\end{Remark}

\subsection*{Доказательство (сведение к неподвижной точке)}

\begin{proof}
	Обозначим $A:=F'_y(0,0)$ --- обратимый оператор. Для каждого фиксированного $x$
	введём вспомогательное отображение $g_x\colon\RR^m\to\RR^m$,
	\[
	g_x(y):=y-A^{-1}F(x,y).
	\]
	Ключевое наблюдение: $y$ --- неподвижная точка $g_x$ (то есть $g_x(y)=y$) $\iff$
	$A^{-1}F(x,y)=0\iff F(x,y)=0$ (так как $A^{-1}$ обратим). Таким образом, решение
	уравнения $F(x,y)=0$ равносильно поиску неподвижной точки $g_x$. Покажем, что при
	малых $x$ отображение $g_x$ --- сжимающее на подходящем шаре, и применим теорему
	Банаха.

	\textbf{Шаг 1 (оценка производной $g_x$).} По $y$ отображение $g_x$
	дифференцируемо, и
	\[
	(g_x)'_y(y)=I_m-A^{-1}F'_y(x,y)=A^{-1}\bigl(A-F'_y(x,y)\bigr)
	=A^{-1}\bigl(F'_y(0,0)-F'_y(x,y)\bigr).
	\]
	Так как $F'_y$ непрерывна в нуле, $\|F'_y(0,0)-F'_y(x,y)\|\to 0$ при $(x,y)\to 0$.
	Поэтому найдётся $\delta>0$ такое, что при $\|x\|<\delta$, $\|y\|<\delta$
	\[
	\|(g_x)'_y(y)\|\le\tfrac12.
	\]

	\textbf{Шаг 2 (сжимаемость).} По теореме о среднем для вектор-функций, для любых
	$y_1,y_2\in B(0,\delta)$
	\[
	\|g_x(y_1)-g_x(y_2)\|
	\le\Bigl(\sup_{\zeta\in[y_1,y_2]}\|(g_x)'_y(\zeta)\|\Bigr)\|y_1-y_2\|
	\le\tfrac12\|y_1-y_2\|,
	\]
	то есть $g_x$ --- сжатие с коэффициентом $\tfrac12$.

	\textbf{Шаг 3 (шар в себя).} Зафиксируем $\theta<\delta$. Так как $F$ непрерывна и
	$F(0,0)=0$, величина $\|F(x,0)\|$ мала при малых $x$; поэтому найдётся
	$\tilde\theta>0$ такое, что при $\|x\|<\tilde\theta$
	\[
	\|g_x(0)\|=\|A^{-1}F(x,0)\|\le\|A^{-1}\|\cdot\|F(x,0)\|\le\tfrac12\theta.
	\]
	Тогда для $y\in\overline{B}(0,\theta)$
	\[
	\|g_x(y)\|\le\|g_x(y)-g_x(0)\|+\|g_x(0)\|\le\tfrac12\|y\|+\tfrac12\theta\le\theta,
	\]
	то есть $g_x\bigl(\overline{B}(0,\theta)\bigr)\subset\overline{B}(0,\theta)$.
	Замкнутый шар $\overline{B}(0,\theta)$ --- полное метрическое пространство.

	\textbf{Шаг 4 (применение теоремы Банаха и непрерывность).} По теореме Банаха при
	каждом $\|x\|<\tilde\theta$ в шаре $\overline{B}(0,\theta)$ существует единственная
	неподвижная точка $y_x$ отображения $g_x$, то есть $F(x,y_x)=0$. Положим
	$\vphi(x):=y_x$ ($U_{x_0}=B(0,\tilde\theta)$, $V_{y_0}=\overline{B}(0,\theta)$);
	единственность неподвижной точки даёт эквивалентность $F(x,y)=0\iff y=\vphi(x)$ в
	$U_{x_0}\times V_{y_0}$ и равенство $\vphi(0)=0$ (так как $F(0,0)=0$).

	Непрерывность $\vphi$ в нуле: выше показано, что для любого сколь угодно малого
	$\theta<\delta$ существует $\tilde\theta>0$, при котором $\|\vphi(x)\|\le\theta$
	для всех $\|x\|<\tilde\theta$. По определению это означает
	$\lim_{x\to 0}\vphi(x)=0=\vphi(0)$.
\end{proof}

\subsection*{Гладкость и производная неявной функции}

\begin{Corollary}{Дифференцируемость $\vphi$}{}
	Если дополнительно $F\in C^1$ в окрестности $(x_0,y_0)$, то $\vphi$ непрерывно
	дифференцируемо в окрестности $x_0$, и его производная вычисляется по формуле
	\[
	\vphi'(x_0)=-\bigl(F'_y(x_0,y_0)\bigr)^{-1}F'_x(x_0,y_0).
	\]
\end{Corollary}

\begin{proof}[Вывод формулы]
	Дифференцируя тождество $F(x,\vphi(x))\equiv 0$ по $x$ (правило цепочки),
	\[
	F'_x(x_0,y_0)+F'_y(x_0,y_0)\,\vphi'(x_0)=0,
	\]
	откуда, пользуясь обратимостью $F'_y(x_0,y_0)$,
	$\vphi'(x_0)=-(F'_y(x_0,y_0))^{-1}F'_x(x_0,y_0)$.
\end{proof}
